We situate Dynamic-BABE in five literatures: classical opinion dynamics, biased assimilation, Social Judgment Theory, co-evolving trust, and algorithmic exposure.

\subsection{Classical Models of Opinion Dynamics}
\citet{DeGroot1974} modelled influence as repeated weighted averaging. In a connected network this process converges to consensus, which cannot explain persistent disagreement. Bounded-confidence models restrict influence to neighbours within a similarity threshold~\citep{Hegselmann2002, Deffuant2000} and can produce multiple clusters. \citet{Flache2017} review these and related formalisms. Bounded confidence typically ignores out-of-range neighbours rather than modelling active repulsion, and classical OD usually lacks an exit channel that changes the communicating population.

\subsection{Cognitive Resistance and Biased Assimilation}
\citet{Lord1979} showed that people can interpret mixed evidence in ways that strengthen prior attitudes---biased assimilation. \citet{Dandekar2013} embedded related biases in network models and linked them to polarization under homophily. \citet{Chen2021} proposed BEBA, where influence weights $w_{ij}(t) = 1 + \beta_i y_i(t) y_j(t)$ depend on entrenchment $\beta_i$ and alignment, so negative influence can emerge endogenously. \citet{Mas2013} discuss bipolarization without explicit distancing. We adapt Chen-style weights to pairwise, multi-issue interactions, map them onto Social Judgment zones~\citep{Sherif1961, Jager2005}, and add co-evolving dyadic trust.

\subsection{Social Judgment Theory and Backfire}
Social Judgment Theory partitions attitude space into latitudes of acceptance, non-commitment, and rejection~\citep{Sherif1961}. Messages in the acceptance latitude tend to pull opinions toward the source; messages in the rejection latitude can produce contrast and backlash. Empirical evidence for attitude \emph{backfire} / boomerang effects is contested: some exposure regimes raise extremity~\citep{Bail2018}, while feed and correction experiments often find null or highly context-dependent attitude change~\citep{Guess2023, Levy2021}. We treat Zone~3 repulsion as a modelled Social Judgment mechanism, not as an empirical universal. \citet{Jager2005} showed that repulsion in networks can create opposing camps. We treat influence weights as determining these zones.

\subsection{Networks and Dyadic Trust}
Opinions and social structure often co-evolve~\citep{Macy2002, Centola2010, Keijzer2018}. Many models allow rewiring toward similar others. We hold the undirected graph fixed and instead let edge trust change with interaction. Trust is asymmetric in gain and loss~\citep{Slovic1993}: it rises slowly after agreement and falls faster after conflict. High trust buffers occasional disagreement; low trust makes further conflict more likely.

\subsection{Algorithmic Interventions, Filtering, and Exposure}
Platforms once hoped that cross-cutting exposure would reduce polarization. \citet{Bail2018} found that exposing Twitter users to opposing-party bots could increase extremity. Feed experiments also show that ranking regimes change on-platform behavior, often more than deep attitudes~\citep{Levy2021, Guess2023}. On the modelling side, algorithmic bias in partner selection can amplify fragmentation and polarization in bounded-confidence settings without modelling exit~\citep{Sirbu2019}; adaptive-network work often rewires ties~\citep{Macy2002, Keijzer2018}. What those neighbours do not jointly deliver---and what we treat as the paper's single distinctive claim---is a suppression filter coupled to Hirschmanian exit so that \emph{retention rises while active-set $\sigma$ stays high or rises}, with $\sigma_{\mathrm{all}}$ showing that the rise is survivorship rather than population warming, plus a \emph{separate} trust-segregation axis on fixed topology. We abstract suppression as a probabilistic filter that intercepts rejection-zone events---distinct from ranking that \emph{promotes} like-minded contact~\citep{Sirbu2019} or constructive bridging. We use ``Polarization Paradox'' for that retention-versus-survivor-dispersion tension without claiming uniqueness of the name. Field studies rarely observe long-run feedback among opinions, trust, and exit; the ABM is built for that gap.
